\usepackage{xeCJK}
\usepackage{geometry}
\geometry{left=2cm,right=2cm,top=2.5cm,bottom=2.5cm}
\setlength{\headheight}{12.7pt}

\usepackage{titlesec}
\titleformat{\section}
  {\normalfont\large\bfseries}
  {\thesection}
  {1em}
  {}
\titlespacing*{\section}
  {0pt}
  {3.5ex plus 1ex minus .2ex}
  {2.3ex plus .2ex}

\usepackage{amsmath}
\usepackage{amssymb}
\usepackage{amsthm}
\usepackage{enumitem}
\setlist[enumerate]{itemsep=0pt,topsep=0pt,parsep=0pt,leftmargin=0pt,itemindent=2.5em}
\setlist[itemize]{itemsep=0pt,topsep=0pt,parsep=0pt,leftmargin=0pt,itemindent=2.5em}
\usepackage{fancyhdr}
\pagestyle{fancy}
\fancyhf{}
\usepackage{graphicx}
\usepackage{hyperref}
\usepackage{indentfirst}
\usepackage{lmodern}


\begin{document}
\setlength{\abovedisplayskip}{1pt}
\setlength{\belowdisplayskip}{1pt}

\section{不交并}

\hypertarget{sec:定义1.1}{}
\noindent\textbf{定义1.1 (不交并)}

给定一族集合$(f(i))_{i\in I}$, 定义它们的\textbf{不交并}
$\displaystyle
\bigsqcup_{i\in I}f(i)\overset{\mathrm{def}}{=}\bigcup_{i\in I}f(i)\times\{i\}$.

\vspace{10pt}

\hypertarget{sec:定义1.2}{}
\noindent\textbf{定义1.2 (嵌入映射)}

给定一族集合$(f(i))_{i\in I}$, 再给定$i\in I$, 则有\textbf{嵌入映射}
$\displaystyle\iota_i:f(i)\to\bigsqcup_{i\in I}f(i),\quad x\mapsto(x,i)$.

显然嵌入映射是单射.

\vspace{10pt}

\hypertarget{sec:定理1.1}{}
\noindent\textbf{定理1.1 (不交并的泛性质)}

给定一族集合$(f(i))_{i\in I}$, 集合$A$和一族映射$(\varphi_i:f(i)\to A)_{i\in I}$,
存在唯一的映射$\displaystyle\varphi:\bigsqcup_{i\in I}f(i)\to A$使得
\[
    \forall i\in I:\,\,\varphi\circ\iota_i=\varphi_i.
\]

\noindent{\kaishu 证明.}

显然$\varphi:(x,i)\mapsto\varphi_i(x)$唯一满足条件. \hfill $\qedsymbol$

\vspace{10pt}

% \hypertarget{sec:定理1.2}{}
% \noindent\textbf{定理1.2}

% 给定一族集合$(f(i))_{i\in I}$和$I$的一个划分$S$, 那么$\displaystyle
% \bigsqcup_{s\in S}\bigsqcup_{i\in s}f(i)\approx\bigsqcup_{i\in I}f(i)$.

% \noindent{\kaishu 证明.}

% 对任意的$\displaystyle((x,i),s)\in\bigsqcup_{s\in S}\bigsqcup_{i\in s}f(i)$,
% 有$s\in S$且$\displaystyle(x,i)\in\bigsqcup_{i\in s}f(i)$,
% 即$i\in s$且$x\in f(i)$. 据此定义
% \[
%     \varphi((x,i),s)=(x,i)\in\bigsqcup_{i\in I}f(i),
% \]
% 即$\displaystyle\varphi:\bigsqcup_{s\in S}\bigsqcup_{i\in s}f(i)
% \to\bigsqcup_{i\in I}f(i)$.

% \vspace{6pt}
% 反之, 对任意的$\displaystyle(x,i)\in\bigsqcup_{i\in I}f(i)$,
% 有$i\in I$且$x\in f(i)$. 于是有唯一的$s_i\in S$使得$i\in s_i$. 据此定义
% \[
%     \phi(x,i)=((x,i),s_i)\in\bigsqcup_{s\in S}\bigsqcup_{i\in s}f(i),
% \]
% 即$\displaystyle\phi:\bigsqcup_{i\in I}f(i)
% \to\bigsqcup_{s\in S}\bigsqcup_{i\in s}f(i)$.

% \vspace{6pt}
% 易证$\varphi,\phi$互为逆映射, 即二者都是双射. \hfill $\qedsymbol$

% \vspace{10pt}

% 有时会使用更简单的二元不交并.

% \vspace{10pt}

% \hypertarget{sec:定义1.3}{}
% \noindent\textbf{定义1.3 (二元不交并)}

% 给定集合$a,b$, 定义它们的\textbf{二元不交并}
% \[
%     a\sqcup b\overset{\mathrm{def}}{=}a\times\{0\}\cup b\times\{1\}.
% \]





% \newpage

\vspace{20pt}

\section{一般 Descartes 积}

\hypertarget{sec:定义2.1}{}
\noindent\textbf{定义2.1 (一般 Descartes 积)}

给定一族集合$(f(i))_{i\in I}$, 定义它们的\textbf{一般Descartes积}
\[
    \prod_{i\in I}f(i)\overset{\mathrm{def}}{=}
    \left\{p:I\to\bigcup_{i\in I}f(i)\,\middle\vert\,
    \forall i\in I:p(i)\in f(i)\right\}.
\]

\vspace{10pt}

\hypertarget{sec:定理2.1}{}
\noindent\textbf{定理2.1}

给定一族集合$(f(i))_{i\in I}$, 如果$f(i)$都非空,
那么$\displaystyle\prod_{i\in I}f(i)$非空.

\noindent{\kaishu 证明.}

取选择函数$\displaystyle F:\{f(i)\mid i\in I\}\to\bigcup_{i\in I}f(i)$.
显然$\displaystyle F\circ f\in\prod_{i\in I}f(i)$. \hfill $\qedsymbol$

\vspace{10pt}

\hypertarget{sec:定义2.2}{}
\noindent\textbf{定义2.2 (投影映射)}

给定一族集合$(f(i))_{i\in I}$, 再给定$i\in I$, 则有\textbf{投影映射}
$\displaystyle\pi_i:\prod_{i\in I}f(i)\to f(i),\quad p\mapsto p(i)$.

显然投影映射是满射.

\vspace{10pt}

\hypertarget{sec:定理2.2}{}
\noindent\textbf{定理2.2 (一般Descartes积的泛性质)}

给定一族集合$(f(i))_{i\in I}$, 集合$A$和一族映射$(\varphi_i:A\to f(i))_{i\in I}$,
存在唯一的映射$\displaystyle\varphi:A\to\prod_{i\in I}f(i)$使得
\[
    \forall i\in I:\,\,\pi_i\circ\varphi=\varphi_i.
\]

\noindent{\kaishu 证明.}

显然$\varphi:a\mapsto(\varphi_i(a))_{i\in I}$唯一满足条件. \hfill $\qedsymbol$

% ``集合的一般构造''这一章就到这里吧. 有缺漏的日后再补.

\end{document}